\section{Méthodologie}
\label{sec:methodology}

\subsection{Vue d'ensemble du pipeline}

Notre pipeline complet comprend quatre étapes : (1) constitution du jeu de données (détection YOLO pour le dataset personnel), (2) prétraitement et augmentation, (3) extraction d'embeddings par un réseau Re-ID, et (4) évaluation par CMC/mAP avec re-ranking optionnel. Pour le clustering (WP7), l'étape 4 est remplacée par un regroupement DBSCAN.

\subsection{Jeux de données}

\paragraph{Market-1501.}
Market-1501 \cite{zheng2015scalable} comprend 32\,668 images de 1\,501 identités capturées par 6 caméras. La partition standard définit 12\,936 images d'entraînement (751 identités), 3\,368 requêtes et 13\,115 images de galerie (750 identités). La résolution des images varie selon la caméra ; toutes sont redimensionnées à $256 \times 128$ px.

\paragraph{Dataset personnel (WP3).}
Pour étudier l'adaptation de domaine, nous avons constitué un jeu de données personnel de 262 images capturées sur le campus de CentraleSupélec. La détection automatique des personnes est réalisée par \textbf{YOLOv8n} \cite{ultralytics2023}, un détecteur léger (\textit{nano}) de la famille Ultralytics YOLO, qui recadre automatiquement chaque personne détectée. Les annotations d'identité (pour l'évaluation) ont été réalisées manuellement. Le dataset est divisé en 131 images d'entraînement et 131 images d'évaluation (requêtes + galerie).

\subsection{Prétraitement et augmentation}

Toutes les images sont redimensionnées à $256 \times 128$ px (ou $224 \times 224$ pour ViT) par interpolation bicubique. Lors de l'entraînement, les augmentations suivantes sont appliquées : retournement horizontal aléatoire ($p=0{,}5$), jitter de couleur (\textit{brightness}, \textit{contrast}, \textit{saturation}, \textit{hue}), et \textit{Random Erasing} -- une occlusion synthétique rectangulaire ($p=0{,}5$). à l'évaluation, seul le redimensionnement et la normalisation ImageNet ($\mu=(0{,}485,\, 0{,}456,\, 0{,}406)$, $\sigma=(0{,}229,\, 0{,}224,\, 0{,}225)$) sont appliqués.

\subsection{Architecture ResNet-50 avec BN-Neck}

Nous utilisons un ResNet-50 \cite{he2016deep} pré-entraîné sur ImageNet, modifié conformément à la baseline BoT \cite{luo2019bag} :

\begin{itemize}
\item \textbf{Last stride = 1} : le stride de la dernière couche convolutive (\texttt{layer4}) est modifié de 2 à 1, doublant la résolution des feature maps ($14 \times 7$ au lieu de $7 \times 4$) et conservant plus d'information spatiale.
\item \textbf{BN-Neck} : apres le pooling global, une couche BatchNorm1d(2048) sans biais produit le vecteur de caractéristiques normalisé $\mathbf{z} \in \mathbb{R}^{2048}$, utilisé pour le calcul des distances à l'inférence. Le classifieur linéaire (sans biais, 2048 $\to$ 751 classes) opère sur $\mathbf{z}$ pendant l'entraînement.
\end{itemize}

Le modèle comporte 25,1M paramètres au total, avec une dimension d'embedding de 2048.

\subsection{Architecture ViT-Base/16 avec BN-Neck}

Nous utilisons ViT-Base/16 \cite{dosovitskiy2021image} pré-entraîné sur ImageNet : l'image $224 \times 224$ est découpée en $196$ patchs de $16 \times 16$ px. Le token \texttt{[CLS]} de dimension 768 est extrait comme représentation globale de l'image ; le même BN-Neck que ResNet-50 lui est appliqué. Le modèle comporte 86,4M paramètres (3,45$\times$ plus que ResNet-50).

\subsection{Stratégies de gel (WP4)}

Pour l'ablation (WP4), trois stratégies de fine-tuning sont étudiées sur ResNet-50 :

\begin{itemize}
\item \textbf{Feature Extraction} : backbone entièrement gelé, seuls BN-Neck et classifieur sont entraînables (1,54M param., 6,2\,\%).
\item \textbf{Partial Fine-Tuning} : \texttt{conv1}, \texttt{layer1} et \texttt{layer2} gelés ; \texttt{layer3}, \texttt{layer4} et tête entraînables (23,6M param., 94,2\,\%).
\item \textbf{Full Fine-Tuning} : tout le réseau entraînable (25,0M param., 100\,\%).
\end{itemize}

\subsection{Protocole d'entraînement}

Les deux modèles sont entraînés avec les hyperparamètres du tableau \ref{tab:training_config}. L'optimiseur Adam est couplé à un schedule cosinus ; la \textit{mixed precision} (AMP) réduit l'empreinte mémoire et accélère l'entraînement. La graine aléatoire est fixée à 42 pour la reproductibilité.

\begin{table}[t]
\centering
\caption{Hyperparamètres d'entraînement.}
\label{tab:training_config}
\resizebox{\linewidth}{!}{%
\begin{tabular}{lcc}
\toprule
Paramètre & ResNet-50 & ViT-Base/16 \\
\midrule
Taux d'apprentissage & $3 \times 10^{-4}$ & $1 \times 10^{-4}$ \\
Batch size & 32 & 24 \\
Optimiseur & Adam & Adam \\
Schedule & Cosinus & Cosinus \\
Epoques & 40 & 40 \\
Résolution input & $256 \times 128$ & $224 \times 224$ \\
Mixed precision & Oui & Oui \\
\bottomrule
\end{tabular}}
\end{table}

\subsection{Re-ranking k-réciproque}

À l'inférence, nous appliquons optionnellement le re-ranking de Zhong \textit{et al.} \cite{zhong2017re} : la matrice de distance cosinus est révisée en tenant compte des voisins k-réciproques ($k_1=20$, $k_2=6$, $\lambda=0{,}3$). Cette procédure améliore significativement la mAP en corrigeant les erreurs de classement local.

\subsection{Descripteurs classiques (WP2)}

Pour la comparaison, trois descripteurs classiques sont utilisés :
\begin{itemize}
\item \textbf{HOG} \cite{dalal2005histograms} : histogramme de gradients orientés sur la totalité de l'image (cellules $8 \times 8$, blocs $16 \times 16$).
\item \textbf{Histogramme de couleur} : histogramme RGB à 3 canaux, 64 bacs par canal.
\item \textbf{SIFT} \cite{lowe2004distinctive} : agrégation de descripteurs locaux de 128 dimensions. Par contrainte de calcul, l'évaluation SIFT est limitée à $n=200$ requêtes et $n=1\,000$ images de galerie (sous-ensemble aléatoire).
\end{itemize}
Le classement est réalisé par recherche exhaustive par distance L2.

\subsection{Clustering DBSCAN (WP7)}

Les embeddings extraits par ResNet-50 sont normalisés puis soumis à DBSCAN \cite{ester1996density} avec distance cosinus. Le paramètre $\varepsilon$ (seuil de densité) est le seul hyperparamètre clé ; nous en étudions la sensibilité systematiquement. Le clustering est évalué par : \textbf{pureté} (proportion d'images correctement assignées à la classe majoritaire de chaque cluster), \textbf{NMI} (\textit{Normalized Mutual Information}), et \textbf{ARI} (\textit{Adjusted Rand Index}).
